\documentclass[11pt]{article} 
\renewcommand{\thempfootnote}{\arabic{mpfootnote}}

\usepackage[hidelinks]{hyperref}

\usepackage[utf8]{inputenc}

\usepackage[T1]{fontenc}
\usepackage[ttdefault=true]{AnonymousPro}

\usepackage{multicol}
\setlength{\columnsep}{1cm}

\usepackage{abc}
\usepackage{graphicx}
\usepackage{musicography}

\usepackage[margin=0.5in]{geometry} 
\geometry{letterpaper}

\usepackage{fancyhdr} 
\pagestyle{fancy} % options: empty , plain , fancy
\renewcommand{\headrulewidth}{0pt} % customise the layout...
\lhead{}\chead{}\rhead{}
\lfoot{}\cfoot{\thepage}\rfoot{}

\begin{document}

\thispagestyle{empty}

\begin{abc}[name=moneymusk,options={-O= +c -F single.fmt}]
X:1
T:Money Musk
T:New England 24-bar version
C:learned from Lissa Schneckenberger
M:C|
K:A
|:d | "A"eAcA "A"eA"D"df | "A"eAcA "E"Bcdf | "A"eAcA "A"cdea | "D"fd"E"Be "A"cAA :|: g |
"A"ae"D"fa "A"eaca | "A"ae"D"db "A"ca"E"Bg | "A"ae"D"fa "A"eaca | "D"fd"E"Be "A"cAA :| B 
|: "A"c2cA "D"d2dA | "A"c2cA "E"B2cB | "A"ABcd "D"efga |1 "E"gefg "A"afed :|2 "E"gefg "A"a4 |]
\end{abc}
% \vskip 0.25in
\noindent\makebox[\textwidth][c]{%
  \begin{minipage}{7.0in}
    \setlength{\parindent}{1.5em}

Money Musk is believed to have been composed by Daniel Dow in 1776 as
a pipe tune (a tune using only the pipe's 9 note range) or a
pipe-inspired Strathspey fiddle
tune. \footnote{\url{https://tunearch.org/wiki/Annotation:Money_Musk\_(1)};
  Goertzen, Chris. George P. Knauff's Virginia Reels and the History
  of American Fiddling, p 48.} \footnote{Strathspey are a type of reel
  that employees the Scots snap rhythm (sixteenth-dotted-eighth-note
  or dotted-eighth-sixteenth-note rhythm) throughout the tune. For an
  audio example of the differences between Strathspey and reel rhythms
  see the YouTube recording of William Coulter, Deby Benton Grosjean,
  Barry Phillips linked on the WOTFA web site.}  `Money Musk' was taken from the
baronial estate, Monymusk House in Aberdeenshire,
Scotland.\footnote{tunearch} The reel Money Musk is played all over
the US, Canada (both French and Anglo Canadians), and the British
Isles, but is particularly popular in Missouri fiddling circles and at
New England contradances.\footnote{tunearch} Money Musk varies greatly
between regions and players using different keys, tune variations, and
overall forms. In New England, Money Musk has its own contradance
steps to complement the 24-bar tune variation (ABC), notably different
to the usual 32-bar tune length (ABAB). I learned the 24-bar New
England version in 2011 from Vermont fiddler, Lissa Schneckenberger,
who stated that it was important to have Money Musk in your back
pocket because there was a high probability that it would be
requested.

Taken from the Contradancers of Hawai'i website:\footnote{\url{http://www.sls.hawaii.edu/bley-vroman/contra/dances/musk.html}}

This version manages the compression by reducing the forward-and-backs
to four-beats--essentially to balances. Note that the two balances are
at different places in the musical phrase. The three-quarter turns in
A2 are at different speeds. The first is quite leisurely, but the
second is fast and must lead directly into the right-and-left. The
balances are sometimes done as forward-and-back balances; sometimes as
step-kick balances. Often these balances are quite percussive. Many
experienced dancers find this a very enjoyable version of the dance:
the differing positions of the balances in the phrase and the
differing speeds of the allemandes add interest. (This was once the
most common version in southwest New Hampshire.\footnote{6 I am not
  from southwest New Hampshire, but my early childhood was spent in
  the Monadnock Region, which overlaps it. Perhaps this is why the
  24-bar version sounds `correct' to me.})

 {\parindent0pt
      \raggedright
      ~\\
A1 Actives allemande right once and a half (8) \\
Go outside, below one couple (4) \\
Balance six (4) \\
A2 Actives allemande right three quarters around, to lines of six across the hall (8) \\
Balance six (4) \\
Actives allemande right three quarters to place (4) \\
B Top two couples (1s and 2s) right and left through (across and back) (16) \\
~\\
 }
 
---Kitty Steetle

    
  \end{minipage}
}


\end{document}
